% !TeX spellcheck = he_IL
%% ! TEX TS-program = xelatexmk
% above: commands appropriate for TeXShop in mac OS
\documentclass[12pt, a4paper]{article}
\usepackage{etoolbox}
\usepackage{bbm}
%\usepackage{pifont}
\usepackage{tikz}
%\input{preamble/UnitPreamble}
%%%%%%%%%%%%%%%%%%%%%%%%%%
%%
%%
%%  Exam Template
%%
%%
%%%%%%%%%%%%%%%%%%%%%%%%%%
\newtoggle{testreview}
\toggletrue{testreview}  
\togglefalse{testreview} % uncomment to make it ready for exam unit
\newtoggle{testsolution}
\toggletrue{testsolution}  
\togglefalse{testsolution} % comment to show solutions
\newcounter{version}
\setcounter{version}{2}  % version of multiple-choice

	





\usepackage{graphicx}

\usepackage{braket}
\usepackage[a4paper, top=3cm,left=2.5cm, right=2.3cm, bottom=2.7cm]{geometry}
\usepackage{amsmath, amssymb, array, multicol, enumitem}
\usepackage{pdfpages}
\usepackage{comment, graphicx}  
\usepackage{xparse}
\DeclareGraphicsExtensions{.png,.jpg}
\usepackage{pgffor}
\usepackage[colorlinks=true]{hyperref} 
\AtBeginDocument{%
	\hypersetup{pdfborder={0 0 0},
		urlcolor=blue}%
}

\usepackage{wasysym}
\usepackage{fontspec}
\usepackage{xunicode}
\usepackage{polyglossia}
\setdefaultlanguage{hebrew}
\setotherlanguage{english}
%\newfontfamily\hebrewfont[Script=Hebrew, Ligatures=TeX, Scale=1.1, WordSpace =1.05]{David}
%\newfontfamily\hebrewfont[Script=Hebrew, Ligatures=TeX, Scale=1.1, WordSpace=1.05]{David  CLM}
\newfontfamily\hebrewfont[Script=Hebrew]{Times New Roman}
%\newfontfamily\hebrewfont[Script=Hebrew]{Ezra SIL}

% Hebrew monospace font (must support Hebrew)
\newfontfamily\hebrewfonttt[Script=Hebrew]{FreeMono}

%http://tex.stackexchange.com/questions/141231/equation-number-enclosing-parentheses-misplaced-with-polyglossiahebrewamsmath
\makeatletter
%\def\maketag@@@#1{\hbox{\m@th\normalfont\LRE{#1}}}
\def\tagform@#1{\maketag@@@{(\ignorespaces#1\unskip)}}
\makeatother

%
% Hebrew fonts are two small 
%
\fontsize{14}{15}

%
% Package fro proper positioning of logo
%
\usepackage[absolute]{textpos}


% http://tex.stackexchange.com/questions/10589/fix-nested-section-numbers-in-rtl-languages-with-polyglossia/
\makeatletter
\renewcommand{\SepMark}[1]{\def\@SepMark{#1}}
\makeatother
\SepMark{.}
\robustify\hebrewnumeral
\robustify\Hebrewnumeral
\robustify\Hebrewnumeralfinal


% from polyglossia docs
%\gappto\captionshebrew{%
\def\urladdrname{{\itshape כתובת רשת}}
\def\emailaddrname{{\itshape כתובת דואל}}
%}

% https://tex.stackexchange.com/questions/176661/polyglossia-arabic-text-english-bibliography
\makeatletter
\let\orig@MakeUppercase \MakeUppercase
\makeatother
\providecommand{\MakeUppercase}[1]{#1}




%%%%%%%%%%%%%%%%%%%%%%%%%%
%%
%%
%%   User Commands
%%
%%
%%%%%%%%%%%%%%%%%%%%%%%%%%
\newcommand{\heb}[1]{\begin{small}\R{#1}\end{small}}
\newcommand{\hebt}[1]{\begin{scriptsize}\R{#1}\end{scriptsize}}
\newcommand{\cc}{\mathbb{C}}
\newcommand{\nn}{\mathbb{N}}
\newcommand{\zz}{\mathbb{Z}}
%\renewcommand{\phi}{\varphi}
\newcommand{\abs}[1]{\left|#1\right|}
\newcommand{\lr}[1]{\left(#1\right)}


\newcommand{\ff}{\mathbb{F}}
\newcommand{\FF}{\mathbb{F}}
\newcommand{\mtn}{m\times n}
\newcommand{\longr}{\longrightarrow}
\newcommand{\longrl}{\longleftrightarrow}
\newcommand{\ol}{\overline}
\newcommand{\ntn}{n\times n}
\newcommand{\rr}{\mathbb{R}}
\newcommand{\QQ}{\mathbb{Q}}
\newcommand{\RR}{\mathbb{R}}
\newcommand{\im}{\mathrm{Im}}

\newcommand{\CC}{\mathbb{C}}
\newcommand{\ZZ}{\mathbb{Z}}
\newcommand{\qq}{\mathbb{Q}}
\DeclareMathOperator{\vol}{vol}

\DeclareMathOperator{\Adj}{Adj}
\DeclareMathOperator{\deter}{det}
\DeclareMathOperator{\Img}{Im}
\DeclareMathOperator{\Ker}{Ker}
\DeclareMathOperator{\rank}{rank}
\DeclareMathOperator{\Span}{Span}
\DeclareMathOperator{\tr}{tr}


\newcommand{\cB}{\mathcal{B}}
\newcommand{\cC}{\mathcal{C}}
\newcommand{\cR}{\mathcal{R}}
\newcommand{\cD}{\mathcal{D}}
\newcommand{\calE}{\mathcal{E}}


%%%% new for complex functions

\newcommand{\Arg}{\mathrm{Arg}}
\newcommand{\wt}{\widetilde}
\newcommand{\wh}{\widehat}
\newcommand{\re}{\mathrm{Re}}
\newcommand{\goes}{\longrightarrow}
\newcommand{\cO}{\mathcal{O}}
\newcommand{\Log}{\mathrm{Log}}


\newcommand{\twopartdef}[4]
{
	\left\{
		\begin{array}{ll}
			#1 & \mbox{   } #2 \\
			#3 & \mbox{   } #4
		\end{array}
	\right.
}

\newcommand{\solutionword}{\textbf{פתרון:}}
\newcommand{\points}[1]{(#1 נק') }
\iftoggle{testsolution}{%
	\newcommand{\sols}[1]{\solutionword$\;\;$ #1$\square$}
	\newcommand{\spacing}{0cm}
	\newcommand{\questionspacing}{2.5cm}
	\newcommand{\extrapages}{}
}{ %
%\newcommand{\spacing}{12cm}
%\newcommand{\halfspacing}{0.5cm}	
\newcommand{\spacing}{24cm}
\newcommand{\questionspacing}{\spacing}
\newcommand{\sols}[1]{}
\newcommand{\extrapages}{%
	\newpage
	\noindent
	דף זה ריק במכוון ומיועד לשימוש הסטודנט/ית.
	\newpage
	\noindent
	דף זה ריק במכוון ומיועד לשימוש הסטודנט/ית.
}
\iftoggle{testreview}{%
\renewcommand{\spacing}{1cm}
%\newcommand{\questionspacing}{\spacing}
\renewcommand{\extrapages}{}	
}{}	
}


\DeclareMathOperator{\row}{Row}

\DeclareMathOperator{\col}{Col}


%%%%%%%%%%%%%%%%%%%%%%%%
%%
%%   BEGIN DOCUMENT
%%
%%%%%%%%%%%%%%%%%%%%%%%%

\begin{document}

\selectlanguage{hebrew}


\newcommand{\semester}{%
	סמסטר א' תשפ״ו
}
\newcommand{\course}{%
 אלגברה לינארית להנדסה 
}
\newcommand{\coursenumber}{9321}	
\newcommand{\lecturers}{%
	מרצים: ד״ר י.מייזל, ד״ר נ.גולקו, ד״ר נ.מאיר, ד״ר נ.איידלשטין, ד״ר מ.פורת, ד״ר מ.זיו-אב, ד״ר מ.ברשדסקי
}

\newcommand{\testdate}{{21.6.26}}
\iftoggle{testsolution}{%
	\newcommand{\testtype}{%
		פתרונות למבחן מועד ג'
	}
}{%
	\newcommand{\testtype}{%
	\boxed{\text{
	מועד ג'
	}	}
}}

\newcommand{\duration}{%
	3 שעות.
}
\newcommand{\helpmaterial}{%
מחשבון פשוט בלבד.
}



\thispagestyle{empty}

\begin{textblock*}{9cm}[0.5,0.5](90mm,3cm)
	\includegraphics[width=5cm]{logo.jpg}
\end{textblock*}	

\noindent

\begin{center}
אוניברסיטת בן-גוריון בנגב - היחידה להוראת קורסי יסוד - %
\smallskip
\semester \\
{\large \bf \course - 214.1.\coursenumber}\\
\smallskip
\lecturers\\
\smallskip
\smallskip
\smallskip
{\Large \testtype }\\
\end{center}

\noindent
\textbf{%
תאריך:
} \testdate \quad 
\textbf{משך הבחינה: 
} \duration \quad
\textbf{
חומר עזר: 
}
\helpmaterial \\
דף הנוסחאות מצורף למבחן. אין להכניס כל חומר עזר נוסף.
\subsection*{חלק א': שאלות רב ברירה}
\begin{itemize}
    \item בחלק א' של המבחן יש 8 שאלות רב-ברירה, ומשקל כל שאלה 5 נקודות.
    \item לכל שאלה יש תשובה אחת נכונה בדיוק. יש לסמן תשובה אחת לכל שאלה \textbf{בטבלה בדף זה}, אין צורך בנימוק.
    \item שימו לב להנחיות מילוי הטבלה בעמוד הבא.
\end{itemize}

\begin{center}
\begin{tabular}{|@{}c@{}||@{}c@{}||@{}c@{}||@{}c@{}||@{}c@{}||@{}c@{}||@{}c@{}||@{}c@{}|}
\hline
1 & 2 &  3 & 4 & 5 & 6 & 7 & 8 \\
\hline
\begin{tabular}{cccc}
\\\textcircled{א}\\\\
\textcircled{ב}\\\\
\Large{\CIRCLE}\\\\
\textcircled{ד}\\\\
\end{tabular}
&
\begin{tabular}{cccc}
\\\textcircled{א}\\\\
\Large{\CIRCLE}\\\\
\textcircled{ג}\\\\
\textcircled{ד}\\\\
\end{tabular}

&
\begin{tabular}{cccc}
\\\Large{\CIRCLE}\\\\
\textcircled{ב}\\\\
\textcircled{ג}\\\\
\textcircled{ד}\\\\
\end{tabular}
&
\begin{tabular}{cccc}
\\\textcircled{א}\\\\
\Large{\CIRCLE}\\\\
\textcircled{ג}\\\\
\textcircled{ד}\\\\
\end{tabular}
&
\begin{tabular}{cccc}
\\\Large{\CIRCLE}\\\\
\textcircled{ב}\\\\
\textcircled{ג}\\\\
\textcircled{ד}\\\\
\end{tabular}
&
\begin{tabular}{cccc}
\\\textcircled{א}\\\\
\textcircled{ב}\\\\
\textcircled{ג}\\\\
\Large{\CIRCLE}\\\\
\end{tabular}
&
\begin{tabular}{cccc}
\\\Large{\CIRCLE}\\\\
\textcircled{ב}\\\\
\textcircled{ג}\\\\
\textcircled{ד}\\\\
\end{tabular}
&
\begin{tabular}{cccc}
\\\Large{\CIRCLE}\\\\
\textcircled{ב}\\\\
\textcircled{ג}\\\\
\textcircled{ד}\\\\
\end{tabular}
\\
\hline 
\begin{tikzpicture}
  \draw[dotted] (0,0) circle (0.8cm);
\end{tikzpicture}
& 
 \begin{tikzpicture}
  \draw[dotted] (0,0) circle (0.8cm);
\end{tikzpicture}
 &
 \begin{tikzpicture}
  \draw[dotted] (0,0) circle (0.8cm);
\end{tikzpicture}
 &
 \begin{tikzpicture}
  \draw[dotted] (0,0) circle (0.8cm);
\end{tikzpicture}
 &
 \begin{tikzpicture}
  \draw[dotted] (0,0) circle (0.8cm);
\end{tikzpicture}
 &
 \begin{tikzpicture}
  \draw[dotted] (0,0) circle (0.8cm);
\end{tikzpicture}
 &
 \begin{tikzpicture}
  \draw[dotted] (0,0) circle (0.8cm);
\end{tikzpicture}
 &
 \begin{tikzpicture}
  \draw[dotted] (0,0) circle (0.8cm);
\end{tikzpicture}
   \\
 \hline
\end{tabular}
\end{center}


\subsection*{חלק ב': שאלות פתוחות}
\begin{itemize}
    \item יש לבחור בדיוק 3 מתוך 4 שאלות פתוחות. משקל כל שאלה 20 נקודות.

    \item יש לסמן את השאלות שנבחרו בטבלה בדף זה. שאלות שלא סומנו לא תבדקנה!
    \item יש לנמק את התשובות בחלק ב' ולכתוב את כל החישובים הרלוונטים.
    \item יש לענות על כל השאלות בחלק זה על טופס המבחן במקום המיועד אחרי כל שאלה.   
    
    \underline{
    ניתן גם לכתוב מצדו השני של הדף. 
    }
    מחברת הטיוטה לא תבדק!
\end{itemize}
\bigskip

 בחלק ב' יש לבדוק את השאלות הבאות. סמנו את התשובות המתאימות --  3 שאלות לכל היותר!


\begin{center}
\setlength{\arrayrulewidth}{0.6pt}
\renewcommand{\arraystretch}{1.4}
\begin{tabular}{|c|c|c|c|}
\hline
\textcircled{1} & \textcircled{2} & \textcircled{3} & \textcircled{4}\\
\hline
\end{tabular}
\end{center}

\begin{center}
\huge{\textbf{בהצלחה!}}
\end{center}
\pagebreak
\subsection*{הוראות למילוי הטבלה}
\begin{itemize}
    \item[$\bullet$] יש להשתמש בעט שחור או כחול. 
    \item[$\bullet$] יש לסמן תשובה אחת בלבד ע"י השחרת העיגול המתאים.
    \item[$\bullet$]
    בשאלה בה לא יושחר אף עיגול,
    הניקוד יהיה 0.
    \item[$\bullet$] בשאלה בה יושחרו יותר מעיגול אחד, הניקוד יהיה 0.
.
    \item[$\bullet$] כדי לתקן סימון, יש לסמן X
    על התשובה הלא נכונה ולהשחיר את התשובה הנכונה.

    \item[$\bullet$]
    \textbf{שימו לב:}
    לא ניתן לבחור סעיף שסומן עליו X
    לכן הקפידו לסמן את תשובתכם כשאתם בטוחים בה.
\end{itemize}

\subsection*{דוגמאות}


\begin{center}
\begin{tabular}{|c||c||c||c||c||c||c||c||c||c|}
\hline
1 & 2 &  3 & 4 & 5 & 6 & 7 & 8 & 9 & 10 \\
\hline
\begin{tabular}{c}
\\\Large{\CIRCLE}\\\\
\ \  \hspace{-8pt}\raisebox{0pt}[0pt][0pt]{\Large\CIRCLE}%
\  \hspace{4pt}\llap{\raisebox{-4pt}[0pt][0pt]{ $\scalebox{3}{$\times$}$}}\hspace{-8pt}
\\\\
\textcircled{ג}\\\\
\textcircled{ד}\\\\
\end{tabular}
&
\begin{tabular}{cccc}
\\\textcircled{א}\\\\
\Large{\CIRCLE}\\\\
\textcircled{ג}\\\\
\textcircled{ד}\\\\
\end{tabular}
&
\begin{tabular}{cccc}
\\
\textcircled{א}
\\\\
\ \  \hspace{-8pt}\raisebox{0pt}[0pt][0pt]{\Large\CIRCLE}%
 \ \hspace{4pt}\llap{\raisebox{-4pt}[0pt][0pt]{ $\scalebox{3}{$\times$}$}}\hspace{-8pt}
\\\\
\Large{\CIRCLE}\\\\
\ \  \hspace{-8pt}\raisebox{0pt}[0pt][0pt]{\Large\CIRCLE}%
 \ \hspace{4pt}\llap{\raisebox{-4pt}[0pt][0pt]{ $\scalebox{3}{$\times$}$}}\hspace{-8pt}
\\\\
\end{tabular}
&
\begin{tabular}{cccc}
\\\textcircled{א}\\\\
\textcircled{ב}\\\\
\textcircled{ג}\\\\
\textcircled{ד}\\\\
\end{tabular}
&
\begin{tabular}{cccc}
\\\hspace{-8pt}\ \  \raisebox{0pt}[0pt][0pt]{\Large\CIRCLE}%
 \ \hspace{4pt}\llap{\raisebox{-4pt}[0pt][0pt]{ $\scalebox{3}{$\times$}$}}\hspace{-8pt}
\\\\
\textcircled{ב}\\\\
\textcircled{ג}\\\\
\textcircled{ד}\\\\
\end{tabular}
&
\begin{tabular}{cccc}
\\\hspace{-8pt}\ \  \raisebox{0pt}[0pt][0pt]{\Large\CIRCLE}%
 \ \hspace{4pt}\llap{\raisebox{-4pt}[0pt][0pt]{ $\scalebox{3}{$\times$}$}}\hspace{-8pt}
\\\\
\hspace{-8pt}\ \  \raisebox{0pt}[0pt][0pt]{\Large\CIRCLE}%
 \ \hspace{4pt}\llap{\raisebox{-4pt}[0pt][0pt]{ $\scalebox{3}{$\times$}$}}\hspace{-8pt}
\\\\
\textcircled{ג}\\\\
\textcircled{ד}\\\\
\end{tabular}
&
\begin{tabular}{cccc}
\\\Large{\CIRCLE}\\\\
\Large{\CIRCLE}\\\\
\textcircled{ג}\\\\
\textcircled{ד}\\\\
\end{tabular}
&
\begin{tabular}{cccc}
\\\Large{\CIRCLE}\\\\
\Large{\CIRCLE}\\\\
\Large{\CIRCLE}\\\\
\textcircled{ד}\\\\
\end{tabular}
&
\begin{tabular}{cccc}
\\\hspace{-8pt}\ \  \raisebox{0pt}[0pt][0pt]{\Large\CIRCLE}%
 \ \hspace{4pt}\llap{\raisebox{-4pt}[0pt][0pt]{ $\scalebox{3}{$\times$}$}}\hspace{-8pt}
\\\\
\Large{\CIRCLE}\\\\
\textcircled{ג}\\\\
\Large{\CIRCLE}\\\\
\end{tabular}
&
\begin{tabular}{cccc}
\\\hspace{-8pt}\ \  \raisebox{0pt}[0pt][0pt]{\Large\CIRCLE}%
 \ \hspace{4pt}\llap{\raisebox{-4pt}[0pt][0pt]{ $\scalebox{3}{$\times$}$}}\hspace{-8pt}
\\\\
\Large{\CIRCLE}
\\\\
\hspace{-8pt}\ \  \raisebox{0pt}[0pt][0pt]{\Large\CIRCLE}%
 \ \hspace{4pt}\llap{\raisebox{-4pt}[0pt][0pt]{ $\scalebox{3}{$\times$}$}}\hspace{-8pt}
\\\\
\Large{\CIRCLE}\\\\
\end{tabular}
\\
\hline
\end{tabular}
\ 

\ 

\begin{itemize}[label=$\bullet$]
    \item שאלות 1,2,3 עם סימונים תקינים עבור תשובות:
    \begin{enumerate}
        \item א
        \item ב
        \item ג
    \end{enumerate}
    \item בשאלות 4,5,6 חסר סימון והניקוד יהיה 0.
    \item בשאלות 7,8,9,10
    יש יותר מסימון אחד והניקוד יהיה 0
\end{itemize}


\end{center}


%%%
%%%
%%%   End of Header and Instructions
%%%
%%%

\pagebreak
%%%%%%%%%%%%%%%%%%%%%%%%%%%%%%%%%%%%%%
%%  Problems
%%%%%%%%%%%%%%%%%%%%%%%%%%%%%%%%%%%%%%


%%%%% שאלה 2 %%%%%%%

\section*{
חלק א' 
(40 נק')
להשחיר את התשובות בטבלה!}
\begin{enumerate}
\item\points{5}
ב-
$\rr^3$
נתון המישור
$$H=\Set{(1,2,0)+t(1,-1,2)+s(2,1,-1) \mid t,s\in \rr}$$
בחרו את הטענה הנכונה:
\begin{enumerate}
\item
המישור
$H$
עובר בראשית הצירים.

\item
$H = \Set{(x,y,z)\in \rr^3\mid x-y-z=-1}$

\item[\fbox{\textcircled{ג}}]
$H=\Set{(x,y,z)\in \rr^3\mid-x+5y+3z=9}$

\item
הנקודה
$(3,1,1)$
שייכת למישור
$H$
\end{enumerate}

\noindent\rule{\textwidth-30pt}{0.4pt}


\item\points{5}
עבור
$a\in\rr$
נתונה המטריצה
$A=
\begin{pmatrix}
a&1&1\\
1&a&1\\
1&1&a
\end{pmatrix}$.

בחרו את הטענה הנכונה:
\begin{enumerate}
\item
עבור
$a=1$
מתקיים
$\dim(N(A))=1$.

\item[\fbox{\textcircled{ב}}]
עבור
$a=-2$
מתקיים
$\rank(A)=2$.

\item
עבור
$a=0$
מתקיים
$\det(A)=0$.

\item
עבור
$a=2$
מתקיים
$\dim(\row(A))=\dim(N(A))$.
\end{enumerate}

\noindent\rule{\textwidth-30pt}{0.4pt}

\item\points{5}
תהי
$T:\rr^2\to\rr^2$
העתקה לינארית המקיימת
$$T(1,1)=(3,1),\qquad T(1,-1)=(1,3)$$
בחרו את הטענה הנכונה:
\begin{enumerate}
\item[\fbox{\textcircled{א}}]
המטריצה המייצגת של
$T$
ביחס לבסיס הסטנדרטי 
$\Set{\begin{pmatrix}
    1 \\ 0
\end{pmatrix},
\begin{pmatrix}
    0 \\ 1
\end{pmatrix}
}$
היא
$\begin{pmatrix}
2&1\\
2&-1
\end{pmatrix}$

\item
המטריצה המייצגת של
$T$
ביחס לבסיס הסטנדרטי
$\Set{\begin{pmatrix}
    1 \\ 0
\end{pmatrix},
\begin{pmatrix}
    0 \\ 1
\end{pmatrix}
}$
היא
$\begin{pmatrix}
2&1\\
1&-1
\end{pmatrix}$

\item
מתקיים
$T(2,0)=(2,2)$.

\item
ההעתקה
$T$
אינה הפיכה.
\end{enumerate}
%%%%% שאלה 2 %%%%%%%
\noindent\rule{\textwidth-30pt}{0.4pt}

\item\points{5}
במרחב
$\rr^3$
נתונים הוקטורים
$$u_1=(1,2,2),\qquad u_2=(2,1,-2),\qquad u_3=(2,-2,1)$$
בחרו את הטענה הנכונה:
\begin{enumerate}
\item
הקבוצה
$\Set{u_1,u_2,u_3}$
היא בסיס אורתונורמלי של
$\rr^3$.

\item[\fbox{\textcircled{ב}}]
הקבוצה
$\Set{u_1,u_2,u_3}$
היא בסיס אורתוגונלי של
$\rr^3$,
אך אינה בסיס אורתונורמלי.

\item
הקבוצה
$\Set{u_1,u_2,u_3}$
אינה בסיס של
$\rr^3$.

\item
קיים
$w\neq (0,0,0)$
כך ש-
$w\perp u_1,u_2,u_3$ (כלומר,
$w$
מאונך לוקטורים
$u_1,u_2,u_3$)
\end{enumerate}

%\noindent\rule{\textwidth-30pt}{0.4pt}
\pagebreak

\item\points{5}
במרחב
$\rr^3$
נתונים הבסיסים הסדורים
$$B=\Set{\begin{pmatrix}
    1 \\ 0 \\1
\end{pmatrix},
\begin{pmatrix}
    0 \\ 1 \\1
\end{pmatrix},
\begin{pmatrix}
    1 \\ 1\\0
\end{pmatrix}},\quad  C=\Set{\begin{pmatrix}
    1 \\ 1 \\1
\end{pmatrix},
\begin{pmatrix}
    1 \\ 1 \\0
\end{pmatrix},
\begin{pmatrix}
    1 \\ 0 \\0
\end{pmatrix}}$$
יהי
$v\in\rr^3$
וקטור המקיים
$$[v]_B=
\begin{pmatrix}
1\\
-1\\
2
\end{pmatrix}$$
בחרו את הטענה הנכונה:
\begin{enumerate}
\begin{multicols}{2}
\item[\fbox{\textcircled{א}}]
$[v]_C=
\begin{pmatrix}
0\\
1\\
2
\end{pmatrix}$

\item
$[v]_C=
\begin{pmatrix}
1\\
-1\\
2
\end{pmatrix}$

\item
$[v]_C=
\begin{pmatrix}
1\\
1\\
0
\end{pmatrix}$

\item
$[v]_C=
\begin{pmatrix}
2\\
-1\\
1
\end{pmatrix}$
\end{multicols}
\end{enumerate}

\noindent\rule{16cm}{0.4pt}
%\pagebreak

%%%%
\item\points{5}
תהיינה
$A,B\in M_{3\times 4}(\rr)$
מטריצות שקולות שורה.
בחרו את הטענה
\underline{
\textbf{שאינה בהכרח}}
נכונה:
\begin{enumerate}
\item
$N(A)=N(B)$

\item
$\rank(A)=\rank(B)$

\item
$\row(A)=\row(B)$

\item[\fbox{\textcircled{ד}}]
$\col(A)=\col(B)$.

\end{enumerate}


\noindent\rule{16cm}{0.4pt}

\item\points{5}
תהי
$A\in M_{3\times 3}(\rr)$
מטריצה הפיכה המקיימת
$\det(A)=-2, \quad \det(A-3I)=5$.

בחרו את הטענה הנכונה:
\begin{enumerate}
\item[\fbox{\textcircled{א}}]
מתקיים
$\det(3A^{-1}-I)=\frac{5}{2}$.

\item
מתקיים
$\det(3A^{-1}-I)=-\frac{5}{2}$.

\item
מתקיים
$\det(3A^{-1}-I)=\frac{5}{54}$.

\item
לא ניתן לקבוע את
$\det(3A^{-1}-I)$
מתוך הנתונים.
\end{enumerate}

\noindent\rule{16cm}{0.4pt}


\item\points{5}
תהי
$A\in M_{3\times 3}(\rr)$
מטריצה שהפולינום האופייני שלה הוא
$$p_A(x)=(x-1)^2(x+2)$$
כמו כן נתון כי
$\rank(A-I)=1$.
בחרו את הטענה הנכונה:
\begin{enumerate}
\item[\fbox{\textcircled{א}}]
המטריצה
$A$
לכסינה מעל
$\rr$.

\item
המטריצה
$A$
אינה לכסינה מעל
$\rr$.

\item
מתקיים
$\dim(N(A+2I))=2$.

\item
מתקיים
$\det(A)=2$.
\end{enumerate}


\pagebreak

\end{enumerate}
%%%
%%%%%%
%%%

 \pagebreak

\section*{חלק ב' (60 נק') - יש לבחור 3 מתוך 4 ולסמן בעמוד הראשון.}
יש לענות במקומות המתאימים בעמודים הבאים.
\begin{enumerate}
\newcommand{\firstQ}{
\item\points{20}
יהי
$a\in\rr$
ותהי
$$A=
\begin{pmatrix}
2&a&0\\
0&2&0\\
0&1&3
\end{pmatrix}$$
\begin{enumerate}
\item\points{5}
מצאו את הפולינום האופייני של
$A$.

\item\points{10}
לכל ערך של
$a\in\rr$
קבעו האם
$A$
לכסינה מעל
$\rr$.

\item\points{5}
עבור
$a=0$
מצאו מטריצה הפיכה
$P$
ומטריצה אלכסונית
$D$
כך ש-
$A=PDP^{-1}$.
\end{enumerate}
}
\firstQ
\noindent\rule{\textwidth-30pt}{0.4pt}

\newcommand{\secondQ}{

\item\points{20}
נסמן
$$B=\left\{
\begin{pmatrix}
1\\
1\\
0
\end{pmatrix},
\begin{pmatrix}
1\\
0\\
1
\end{pmatrix},
\begin{pmatrix}
0\\
1\\
1
\end{pmatrix}
\right\}$$
תהי
$T:\rr^3\to\rr^3$
העתקה לינארית המקיימת
$$T\begin{pmatrix}
1\\
1\\
0
\end{pmatrix}
=
\begin{pmatrix}
3\\
2\\
2
\end{pmatrix},
\qquad
T\begin{pmatrix}
1\\
0\\
1
\end{pmatrix}
=
\begin{pmatrix}
2\\
1\\
2
\end{pmatrix},
\qquad
T\begin{pmatrix}
0\\
1\\
1
\end{pmatrix}
=
\begin{pmatrix}
3\\
1\\
4
\end{pmatrix}$$

\begin{enumerate}
\item\points{8}
מצאו נוסחה עבור
$$T\begin{pmatrix}
a\\
b\\
c
\end{pmatrix}$$

\item\points{8}
מצאו בסיס אורתונורמלי עבור
$\ker(T)$
ובסיס אורתונורמלי עבור
$\im(T)$.

\item\points{4}
האם קיימים
$v_1,v_2\in\rr^3$
כך ש-
$v_1\neq v_2$
וגם
$T(v_1)=T(v_2)$?
נמקו.
\end{enumerate}

}
\secondQ

%\noindent\rule{\textwidth-30pt}{0.4pt}
\pagebreak

\newcommand{\thirdQ}{

\item\points{20}
יהי
$a\in\rr$
ותהי
$$A=
\begin{pmatrix}
1&2&1\\
2&5&3\\
1&3&a
\end{pmatrix}.$$

\begin{enumerate}
\item\points{12}
לכל ערך של
$a\in\rr$
מצאו בסיס עבור
$\row(A)$,
בסיס עבור
$\col(A)$
ובסיס עבור
$N(A)$.

\item\points{4}
לכל ערך של
$a\in\rr$
מצאו בסיס עבור
$\row(A)\cap N(A)$.

\item\points{4}
לכל ערך של
$a\in\rr$
מצאו בסיס עבור
$\col(A)+N(A)$.
\end{enumerate}

}
\thirdQ

\noindent\rule{\textwidth-30pt}{0.4pt}
\newcommand{\forthQ}{

\item\points{20}
ענו על שני הסעיפים הבאים.

\begin{enumerate}
\item\points{10}
תהי
$T:\rr^n\to\rr^n$
העתקה לינארית המקיימת
$T( T(v))={0}$
לכל
$v\in \rr^n$.

הוכיחו כי
$\im(T)\subseteq \ker(T)$
והסיקו כי
$\dim(\im(T))\leq \frac{n}{2}$

\bigskip

\item\points{10}
תהי
$T:V\to W$
העתקה לינארית חד-חד ערכית,

ויהיו
$v_1,v_2,v_3\in V$
וקטורים בלתי תלויים לינארית (בת"ל).

הוכיחו כי
$T(v_1),T(v_2),T(v_3)$
בלתי תלויים לינארית (בת"ל).
\end{enumerate}

}\forthQ
\pagebreak

\setcounter{enumi}{0}
    \firstQ
\noindent\rule{\textwidth-30pt}{0.4pt}
\underline{פיתרון:}

\begin{enumerate}
\item הפולינום האופייני הוא
\begin{align*}
p_A(x)
&=\det(xI-A)\\
&=\det\begin{pmatrix}
x-2&-a&0\\
0&x-2&0\\
0&-1&x-3
\end{pmatrix}\\
&=(x-2)^2(x-3).
\end{align*}

\item הערכים העצמיים הם
$2$
עם ריבוי אלגברי
$2$
ו-
$3$
עם ריבוי אלגברי
$1$.
עבור
$\lambda=2$
מקבלים
\begin{align*}
A-2I=
\begin{pmatrix}
0&a&0\\
0&0&0\\
0&1&1
\end{pmatrix}.
\end{align*}
אם
$a\neq0$
אז
$y=0$
ו-
$z=0$,
ולכן
$\dim N(A-2I)=1$.
במקרה זה אין מספיק וקטורים עצמיים, ולכן
$A$
אינה לכסינה.

אם
$a=0$
אז התנאי היחיד הוא
$y+z=0$,
ולכן
$\dim N(A-2I)=2$.
בנוסף ל-
$\lambda=3$
יש מרחב עצמי מממד
$1$,
ולכן במקרה זה
$A$
לכסינה.

לסיכום:
$$A\text{ לכסינה מעל }\rr \iff a=0.$$

\item עבור
$a=0$
אפשר לבחור וקטורים עצמיים
\begin{align*}
\lambda=2:&\qquad
v_1=\begin{pmatrix}1\\0\\0\end{pmatrix},\quad
v_2=\begin{pmatrix}0\\-1\\1\end{pmatrix},\\
\lambda=3:&\qquad
v_3=\begin{pmatrix}0\\0\\1\end{pmatrix}.
\end{align*}
לכן אפשר לקחת
\begin{align*}
P=
\begin{pmatrix}
1&0&0\\
0&-1&0\\
0&1&1
\end{pmatrix},
\qquad
D=
\begin{pmatrix}
2&0&0\\
0&2&0\\
0&0&3
\end{pmatrix},
\end{align*}
ואז
$A=PDP^{-1}$.
\end{enumerate}

    \pagebreak

    \secondQ
    \noindent\rule{\textwidth-30pt}{0.4pt}
    \underline{פיתרון:}

\begin{enumerate}
\item
נסמן
$$
v_1=(1,1,0),
\qquad
v_2=(1,0,1),
\qquad
v_3=(0,1,1).
$$
צריך לכתוב את
$(a,b,c)$
כצירוף לינארי של איברי
$B$.
נפתור
$$
(a,b,c)
=
\alpha(1,1,0)+\beta(1,0,1)+\gamma(0,1,1).
$$
כלומר
$$
\alpha+\beta=a,
\qquad
\alpha+\gamma=b,
\qquad
\beta+\gamma=c.
$$
מכאן
$$
\alpha=\frac{a+b-c}{2},
\qquad
\beta=\frac{a-b+c}{2},
\qquad
\gamma=\frac{-a+b+c}{2}.
$$
לכן
$$
T(a,b,c)
=
\alpha(3,2,2)+\beta(2,1,2)+\gamma(3,1,4).
$$
נציב ונקבל
$$
T(a,b,c)
=
(a+2b+c,\ a+b,\ 2b+2c).
$$

\item
נמצא בסיס אורתונורמלי עבור
$\ker(T)$.
נפתור
$$
T(a,b,c)=(0,0,0).
$$
כלומר
$$
a+2b+c=0,
\qquad
a+b=0,
\qquad
2b+2c=0.
$$
מהמשוואה השנייה נקבל
$a=-b$,
ומהמשוואה השלישית נקבל
$c=-b$.
לכן
$$
(a,b,c)=(-b,b,-b)=t(1,-1,1),
$$
ומכאן
$$
\ker(T)=\Span\{(1,-1,1)\}.
$$
ננרמל:
$$
\|(1,-1,1)\|=\sqrt{3}.
$$
לכן בסיס אורתונורמלי עבור
$\ker(T)$
הוא
$$
\boxed{
\left\{
\frac{1}{\sqrt{3}}(1,-1,1)
\right\}
}.
$$

כעת נמצא בסיס אורתונורמלי עבור
$\im(T)$.
המטריצה המייצגת של
$T$
ביחס לבסיס הסטנדרטי היא
$$
A=
\begin{pmatrix}
1&2&1\\
1&1&0\\
0&2&2
\end{pmatrix}.
$$
לכן
$\im(T)=\col(A)$.
ניקח בסיס עבור מרחב העמודות:
$$
v_1=(1,1,0),
\qquad
v_2=(2,1,2).
$$
נפעיל את תהליך גראם-שמיט לפי הנוסחה, וננרמל רק בסוף.
נגדיר
$$
w_1=v_1=(1,1,0).
$$
כעת
$$
w_2
=
v_2-\frac{v_2\cdot w_1}{w_1\cdot w_1}w_1.
$$
נחשב
$$
v_2\cdot w_1=3,
\qquad
w_1\cdot w_1=2.
$$
ולכן
$$
w_2
=
(2,1,2)-\frac{3}{2}(1,1,0)
=
\left(\frac{1}{2},-\frac{1}{2},2\right).
$$
אפשר להכפיל ב-
$2$
ולקבל בסיס אורתוגונלי פשוט יותר:
$$
\{(1,1,0),(1,-1,4)\}.
$$
כעת ננרמל:
$$
\|(1,1,0)\|=\sqrt{2},
\qquad
\|(1,-1,4)\|=\sqrt{18}=3\sqrt{2}.
$$
לכן בסיס אורתונורמלי עבור
$\im(T)$
הוא
$$
\boxed{
\left\{
\frac{1}{\sqrt{2}}(1,1,0),
\frac{1}{3\sqrt{2}}(1,-1,4)
\right\}
}.
$$

\item
כן.
מהסעיף הקודם
$$
\ker(T)=\Span\{(1,-1,1)\},
$$
ולכן
$\ker(T)\neq\{0\}$.
למשל ניקח
$$
v_1=(0,0,0),
\qquad
v_2=(1,-1,1).
$$
אז
$v_1\neq v_2$
אבל
$$
T(v_1)=T(v_2)=(0,0,0).
$$
\end{enumerate}

    \pagebreak

    \thirdQ
    \noindent\rule{\textwidth-30pt}{0.4pt}
    \underline{פיתרון:}

\begin{enumerate}
\item
נחשב דטרמיננטה:
$$
\det(A)
=
\begin{vmatrix}
1&2&1\\
2&5&3\\
1&3&a
\end{vmatrix}
=
1(5a-9)-2(2a-3)+(6-5)
=
a-2.
$$
לכן אם
$a\neq 2$
אז
$A$
הפיכה, ולכן
$$
\row(A)=\rr^3,
\qquad
\col(A)=\rr^3,
\qquad
N(A)=\{0\}.
$$
במקרה זה בסיסים מתאימים הם
$$
\boxed{
\row(A):\ \{(1,0,0),(0,1,0),(0,0,1)\}
}
$$
$$
\boxed{
\col(A):\ \{(1,0,0),(0,1,0),(0,0,1)\}
}
$$
$$
\boxed{
N(A):\ \emptyset
}.
$$

נניח כעת
$a=2$.
אז
$$
A=
\begin{pmatrix}
1&2&1\\
2&5&3\\
1&3&2
\end{pmatrix}.
$$
נדרג:
$$
\begin{pmatrix}
1&2&1\\
2&5&3\\
1&3&2
\end{pmatrix}
\longrightarrow
\begin{pmatrix}
1&2&1\\
0&1&1\\
0&1&1
\end{pmatrix}
\longrightarrow
\begin{pmatrix}
1&2&1\\
0&1&1\\
0&0&0
\end{pmatrix}
\longrightarrow
\begin{pmatrix}
1&0&-1\\
0&1&1\\
0&0&0
\end{pmatrix}.
$$
לכן בסיס עבור מרחב השורות הוא
$$
\boxed{
\row(A)=\Span\{(1,0,-1),(0,1,1)\}
}.
$$
העמודות המובילות הן העמודות הראשונה והשנייה, ולכן בסיס עבור מרחב העמודות הוא
$$
\boxed{
\col(A)=\Span\{(1,2,1),(2,5,3)\}
}.
$$
כדי למצוא את הגרעין נפתור
$$
x-z=0,
\qquad
y+z=0.
$$
לכן
$$
(x,y,z)=(t,-t,t)=t(1,-1,1),
$$
ומקבלים
$$
\boxed{
N(A)=\Span\{(1,-1,1)\}
}.
$$

\item
אם
$a\neq 2$
אז
$A$
הפיכה, ולכן
$N(A)=\{0\}$.
מכאן
$$
\row(A)\cap N(A)=\{0\}.
$$

נניח כעת
$a=2$.
מסעיף קודם קיבלנו
$$
N(A)=\Span\{(1,-1,1)\}.
$$
בנוסף,
$A$
סימטרית, ולכן
$$
\row(A)=\col(A).
$$
נראה כי
$(1,-1,1)\notin\col(A)$.
עבור
$a=2$
עמודות
$A$
הן
$$
c_1=(1,2,1),
\qquad
c_2=(2,5,3),
\qquad
c_3=(1,3,2).
$$
אם
$(1,-1,1)\in\col(A)$,
אז קיימים
$\alpha,\beta,\gamma\in\rr$
כך ש-
$$
\alpha c_1+\beta c_2+\gamma c_3=(1,-1,1).
$$
כלומר
$$
\begin{cases}
\alpha+2\beta+\gamma=1\\
2\alpha+5\beta+3\gamma=-1\\
\alpha+3\beta+2\gamma=1
\end{cases}
$$
מהמשוואה השלישית פחות הראשונה נקבל
$$
\beta+\gamma=0,
$$
ומהמשוואה השנייה פחות פעמיים הראשונה נקבל
$$
\beta+\gamma=-3.
$$
זו סתירה, ולכן
$$
(1,-1,1)\notin\col(A)=\row(A).
$$
מאחר ש-
$N(A)=\Span\{(1,-1,1)\}$,
נקבל
$$
\row(A)\cap N(A)=\{0\}.
$$
לכן לכל
$a\in\rr$
בסיס עבור
$\row(A)\cap N(A)$
הוא
$$
\boxed{\emptyset}.
$$

\item
אם
$a\neq 2$
אז
$A$
הפיכה, ולכן
$$
\col(A)=\rr^3,
\qquad
N(A)=\{0\}.
$$
לכן
$$
\col(A)+N(A)=\rr^3,
$$
ובסיס אפשרי הוא
$$
\boxed{
\{(1,0,0),(0,1,0),(0,0,1)\}
}.
$$

נניח כעת
$a=2$.
מסעיף קודם קיבלנו
$$
\dim(\col(A))=2,
\qquad
\dim(N(A))=1.
$$
בנוסף, מסעיף ב' קיבלנו
$$
\col(A)\cap N(A)=\{0\},
$$
כי
$\row(A)=\col(A)$.
לפי משפט המימדים,
$$
\dim(\col(A)+N(A))
=
\dim(\col(A))+\dim(N(A))-\dim(\col(A)\cap N(A))
=
2+1-0
=
3.
$$
לכן
$$
\col(A)+N(A)=\rr^3.
$$
בסיס מתאים הוא
$$
\boxed{
\{(1,2,1),(2,5,3),(1,-1,1)\}
}.
$$
\end{enumerate}
    \pagebreak

    \forthQ
    \noindent\rule{\textwidth-30pt}{0.4pt}
    \underline{פיתרון:}

\begin{enumerate}
\item יהי
$w\in\im(T)$.
אז קיים
$v\in\rr^n$
כך ש-
$w=T(v)$.
לכן
\begin{align*}
T(w)=T(T(v))=0.
\end{align*}
כלומר
$w\in\ker(T)$,
ולכן
$\im(T)\subseteq\ker(T)$.

כעת לפי משפט המימד להעתקות לינאריות,
\begin{align*}
\dim(\im(T))+\dim(\ker(T))=n.
\end{align*}
מאחר ש-
$\im(T)\subseteq\ker(T)$,
נקבל
$\dim(\im(T))\leq\dim(\ker(T))$.
לכן
\begin{align*}
2\dim(\im(T))
\leq
\dim(\im(T))+\dim(\ker(T))
=n,
\end{align*}
ומכאן
\begin{align*}
\dim(\im(T))\leq\frac{n}{2}.
\end{align*}

\item נניח כי
\begin{align*}
\alpha_1T(v_1)+\alpha_2T(v_2)+\alpha_3T(v_3)=0.
\end{align*}
מלינאריות של
$T$
נקבל
\begin{align*}
T(\alpha_1v_1+\alpha_2v_2+\alpha_3v_3)=0.
\end{align*}
מכיוון ש-
$T$
חד-חד ערכית, מתקיים
$\ker(T)=\{0\}$.
לכן
\begin{align*}
\alpha_1v_1+\alpha_2v_2+\alpha_3v_3=0.
\end{align*}
מאחר ש-
$v_1,v_2,v_3$
בלתי תלויים לינארית, נקבל
\begin{align*}
\alpha_1=\alpha_2=\alpha_3=0.
\end{align*}
לכן
$T(v_1),T(v_2),T(v_3)$
בלתי תלויים לינארית.
\end{enumerate}

    \pagebreak
\end{enumerate}

    \pagebreak 
    
    עמוד נוסף לכתיבת המשך פיתרונות
    
    \pagebreak
    עמוד נוסף לכתיבת המשך פיתרונות

 \pagebreak

 % --- דף הנוסחאות ---

\section*{דף הנוסחאות}

\begin{enumerate}
\item $M_{m\times n}(F)$ זו קבוצת המטריצות בגודל $m\times n$
($m$ שורות, $n$ עמודות)
עם רכיבים מ-$F$.

\noindent\rule{\textwidth-30pt}{0.4pt}

\item לכל מטריצה $A$, נסמן על ידי $(A)_{i,j}$ את האיבר של $A$ הנמצא בשורה ה-$i$ ובעמודה ה-$j$.

\noindent\rule{\textwidth-30pt}{0.4pt}

\item פיתוח דטרמיננטה של $A\in M_{n\times n}(F)$ לפי השורה ה-$i$ או לפי העמודה ה-$j$:
\[
\det(A)=\sum_{i=1}^{n}(-1)^{i+j}a_{ij}\det\!\big(M_{ij}(A)\big)
=\sum_{j=1}^{n}(-1)^{i+j}a_{ij}\det\!\big(M_{ij}(A)\big)
\]
כאן $M_{ij}(A)$ זה המינור ה-$i,j$ של המטריצה $A$.

\noindent\rule{\textwidth-30pt}{0.4pt}

\item אם $T:V\to W$ העתקה לינארית, אם $B=\{v_1,\ldots,v_n\}$ בסיס סדור של $V$ ואם $C$ בסיס סדור של $W$,
אז המטריצה המייצגת של $T$ ביחס לבסיסים $B$ ו-$C$ מוגדרת על ידי
\[
[T]^B_C=
\begin{pmatrix}
\vline & & \vline\\
[T(v_1)]_C & \cdots & [T(v_n)]_C\\
\vline & & \vline
\end{pmatrix}.
\]

\noindent\rule{\textwidth-30pt}{0.4pt}

\item לכל $u=(a_1,\ldots,a_n),\; w=(b_1,\ldots,b_n)\in\mathbb{R}^n$

נסמן את המכפלה הפנימית הסטנדרטית ב-$\mathbb{R}^n$ (המכפלה הסקלרית):
\[
 u\cdot w=a_1\cdot b_1+\cdots+a_n\cdot b_n.
\]

\noindent\rule{\textwidth-30pt}{0.4pt}

% --- דף הנוסחאות ---


\item תהליך גראם-שמיט (גרסה א):
אם $W$ תת-מרחב של $\rr^n$, ואם $\{v_1,\ldots,v_m\}$ בסיס של $W$

אז $\{w_1,\ldots,w_m\}$ בסיס אורתוגונלי של $W$ כאשר
\[
w_1=v_1,\qquad
w_k=v_k-\sum_{i=1}^{k-1}\frac{ v_k\cdot w_i}{ w_i\cdot w_i}\,w_i
\quad (k=2,\ldots,m).
\]
\item תהליך גראם-שמיט (גרסה ב):
בהנתן סדרה בת"ל
$\{v_1,\ldots,v_k\}\subseteq \rr^n$, נבנה סדרה אורתונורמלית באופן הבא:
\[
w_1=\frac{v_1}{\|v_1\|},\qquad
w_i'=v_i-\sum_{l=1}^{i-1} (v_i\cdot w_l)\,w_l,\qquad w_i = \frac{w_i'}{\|w_i'\|}
\]

\noindent\rule{\textwidth-30pt}{0.4pt}

\item השורשים של המשוואה הריבועית
$
ax^2+bx+c=0\ (a\neq 0)
$
נתונים על ידי
\[
x_{1,2}=\frac{-b\pm\sqrt{b^2-4ac}}{2a}.
\]

\noindent\rule{\textwidth-30pt}{0.4pt}

\item מטריצה $A\in M_{n\times n}(\ff)$ דומה למטריצה $B\in M_{n\times n}(\ff)$ 

אם
קיימת 
מטריצה הפיכה
$P\in M_{n\times n}(\ff)$ כך ש-
$A=PBP^{-1}$.
\end{enumerate}

\end{document}